\section{Detailed Design}
\label{sec:design}

This section walks the four flows that carry most of the system's behaviour.
Each is illustrated with the relevant sequence diagram and a short
discussion of failure modes.

\subsection{Voting flow}
\label{sub:vote-flow}

The vote flow is the most security-critical because it is the only
endpoint open to anonymous traffic. We layer four defences:
\textbf{(a)} schema validation, \textbf{(b)} per-IP rate limiting,
\textbf{(c)} an HMAC \emph{proof-of-presence} token issued at \texttt{/init}
and verified at \texttt{/}, and \textbf{(d)} an in-memory anomaly scorer.
Figure~\ref{fig:seq-vote} shows the full sequence.

\fig{seq-vote}{0.96}{Sequence diagram for the two-step vote flow.}

\paragraph{Why two steps?} A single-endpoint design (just \texttt{POST /votes})
is trivially scriptable: a curl loop with a UUID per iteration would pass
every check we do. Splitting it forces a stateful round-trip; the
attacker must call \texttt{/init}, store the token, and post it back inside
the 5-minute TTL; which is enough friction that the rate limit and
anomaly detector dominate the cost-benefit curve.

\paragraph{What the HMAC binds.} The token's payload is
\texttt{browserToken|issuedAt|nonce}. Tampering with any of those three
invalidates the signature. Reuse across browser tokens fails because
\texttt{verifyVoteToken} compares the embedded browser token to the one
in the cast payload using \texttt{timingSafeEqual} (constant-time, to
defend against the unlikely-but-cheap timing attack).

\paragraph{Anomaly scoring.} Figure~\ref{fig:flow-anomaly} shows the
scoring rules. The thresholds were tuned on synthetic load to flag
$\geq 95\%$ of scripted bursts at severity \texttt{high} while keeping
the false-positive rate at zero on a 100-vote organic baseline.

\fig{flow-anomaly}{0.85}{Vote anomaly scoring rules.}

\paragraph{Failure modes.}
\begin{itemize}[leftmargin=*]
  \item Supabase write fails $\rightarrow$ controller returns 500, the
        client surfaces a generic ``Could not record vote'' message; the
        caller can retry.
  \item Supabase \emph{read} (the duplicate check) fails $\rightarrow$ we
        treat absence as ``not yet voted''; worst case is one duplicate
        vote slipping through, which is then flagged by the anomaly
        detector on its second arrival.
  \item Anomaly persistence (\texttt{vote\_flags} insert) fails $\rightarrow$
        logged at WARN, vote still counts. Integrity dashboard misses the
        flag for that one vote; we accept this as the lesser failure.
\end{itemize}

\subsection{Grading flow}
\label{sub:grade-flow}

Figure~\ref{fig:seq-grade} shows the auto-save plus submit flow.

\fig{seq-grade}{0.95}{Sequence diagram for grading.}

The grading interface debounces the user's keystrokes at 800\,ms and posts
to \texttt{PUT /grades/:id/draft}; the controller upserts a row with
\texttt{status='in\_progress'}. When the judge clicks Submit, the
controller verifies \texttt{status != 'submitted'}, validates that all
five scores are filled, and upserts with \texttt{status='submitted'}.
After a successful submit we (i) write a row to \texttt{audit\_logs},
(ii) invalidate the cached \texttt{winners:} keys, and (iii) broadcast
\texttt{grade.submit} on the event bus.

The state machine of a grade row is shown in
Figure~\ref{fig:state-grade}. Note the absence of an arrow back from
\texttt{Submitted}: there is no controller path that walks the row back
to \texttt{InProgress}. An admin who must reset a grade does so via the
Supabase SQL editor, and that out-of-band action would itself be visible
to anyone replaying the audit log because the next \texttt{grade.submit}
would carry a different \texttt{previousStatus} in its metadata.

\fig{state-grade}{0.65}{Grade status machine.}

\subsection{Real-time session sync}
\label{sub:sse-flow}

The admin presses ``make event public'' in one browser; every public
voter's home page updates within 200\,ms. Figure~\ref{fig:seq-sse}
shows how.

\fig{seq-sse}{0.96}{Server-Sent Events session push.}

The browser opens \texttt{GET /api/session/stream}. The server registers
the connected response in a process-local \texttt{Set} (the
\emph{event bus}). On any session-affecting write, the controller writes
the durable change to Supabase, then calls
\texttt{eventBus.broadcast('session.update', shapedSession)} which
\texttt{res.write}s to every member of the set.

A heartbeat \texttt{: ping} (an SSE comment line) goes out every 25\,s on
all sockets. Render's idle proxy timeout is 60\,s, so this keeps the
socket alive indefinitely. If the socket nevertheless drops, the
client's \texttt{useEventStream} watchdog fires after \texttt{staleMs}
(default 45\,s) of silence and reconnects.

\paragraph{Why a watchdog and a polling fallback?} The native
\texttt{EventSource} \emph{does} reconnect on errors, but on Render's
free tier a clean idle close (FIN, no error event) sometimes occurs and
\texttt{onerror} never fires. The watchdog is a defence against that
specific known-known. The 60-second polling fallback is a defence
against unknown unknowns.

\subsection{AI insight generation}
\label{sub:ai-flow}

Figure~\ref{fig:seq-ai} shows the cache-first read of an insight.

\fig{seq-ai}{0.96}{AI insight retrieval and generation.}

The first request for a project hits the
\texttt{ai\_insights} table, finds nothing, fetches the project row,
calls Claude with a cached system prompt, parses the JSON out of the
response, persists it, and returns. Subsequent requests in the same
session find the row and skip the API call entirely.

\paragraph{Prompt caching.} The system prompt block is declared with
\texttt{cache\_control: \{type: 'ephemeral'\}}. Anthropic caches a hash
of the cached prefix for five minutes. The first call writes the cache
(at $1.25\times$ the standard input rate); every subsequent call within
that window reads the cache (at $0.1\times$ the standard rate). For 9
featured projects pre-warmed in a tight loop this means the system
prompt is paid for once, not nine times.

\paragraph{The deterministic stub.} If \texttt{ANTHROPIC\_API\_KEY} is
absent (or the call fails), \texttt{stubInsights} returns a payload of
the same shape, derived from the project description. The UI never
distinguishes \emph{real} from \emph{stub}; only a small ``offline mode''
chip appears for transparency.

\subsection{Auth flow}

The auth flow is the most boring of the four but worth documenting.
Figure~\ref{fig:flow-auth} captures both the login path and the
\texttt{protect} middleware on subsequent requests.

\fig{flow-auth}{0.92}{Authentication and protected-request flow.}

\subsection{Caching strategy}

The \texttt{cache} module (Listing~\ref{lst:cache}) exposes
\texttt{getOrSet}, \texttt{invalidate}, and \texttt{stats}. The two
disciplines we follow are: \textbf{(i)} every read controller wraps its
DB call in \texttt{getOrSet}, \textbf{(ii)} every write controller calls
\texttt{invalidate} on the prefix(es) it just made stale.

The cache key is a structured string with a colon-delimited
prefix-key form (e.g.\ \texttt{projects:project:All::p1:l50}); this lets
\texttt{invalidate} use a simple \texttt{startsWith} match.

\subsection{Anomaly detector internals}

The detector's data structure is a single in-memory ring of recent vote
events with the shape \texttt{\{ip, ua, time, browserToken\}}. On every
arrival we trim entries older than 5 minutes, then compute six
contributions (Figure~\ref{fig:flow-anomaly}). The total maps to a
severity enum.

The intentional simplicity of these rules is part of the design. They
are explainable to a non-engineer (``three votes from one IP in a
minute? flag''), and explainability matters when an admin has to
override the system on the day of a real event.
